\documentclass[../main]{subfiles}
\begin{document}
\ifSubfilesClassLoaded{\setcounter{page}{4}}{}
\section{Comunismo e organizzazione della vita sociale}
\label{sec:07_comunismo_gestione}

Nella fase di transizione verso una società comunista, la pianificazione e la gestione del tempo libero richiedono l’intervento dello Stato. Tale intervento, tuttavia, si rende necessario solo in quanto la produzione del tempo libero coincide con la produzione e la distribuzione dei beni essenziali per la collettività, e soltanto fino a quando questa produzione non sarà organizzata secondo principi comunisti. Il tempo libero, in quanto tale, e quindi intrinsecamente legato all’ideale comunista, sfugge a qualsiasi forma di controllo statale esterno.

Anche nella società capitalista odierna, il tempo libero rappresenta uno spazio in cui l’autorità dello Stato non si esercita. In questo ambito, quando si praticano attività libere e ispirate a principi comunitari, la gestione è affidata a tutti i membri di una libera associazione. Con l’incremento delle attività improntate a valori comunitari, il ruolo dello Stato tende a ridursi.

Incentivare la produzione di attività ricreative e, di conseguenza, di forme di attività collettiva libera, significa compiere un atto volto a minare le condizioni di esistenza delle classi durante il periodo di transizione, e allo stesso tempo un atto di autodistruzione dello Stato. Lo Stato è, in sostanza, un apparato che tutela gli interessi della classe dominante, esercitando un’amministrazione a beneficio di tale classe, e questa sua funzione primaria viene meno con la dissoluzione delle classi, inclusa la classe dominante.

Di conseguenza, il ruolo dello Stato nel periodo di transizione si riduce a garantire la limitazione (l’eliminazione) del carattere mercantile della produzione e a farsi carico della creazione delle condizioni per un’attività libera e comunitaria dei membri della società.

La prima di queste sfide si pone con particolare evidenza di fronte agli Stati capitalisti. Il capitalismo, per sua natura, è fonte inesauribile di conflitti, e la logica bellica esige un sistema di distribuzione centralizzata dei beni destinati al consumo individuale. Uno Stato che si dimostrasse incapace di assolvere a questo compito sarebbe condannato all'estinzione, sia nella sua struttura istituzionale, sia insieme alla popolazione che governa. La realizzazione di un tale sistema comporta, in molteplici aspetti, la necessità di subordinare gli interessi di singoli capitalisti all'interesse generale della classe capitalista, un obiettivo che non sempre può essere conseguito appieno e con il grado di efficacia desiderato.

In tempo di guerra, l’erogazione di beni di prima necessità e di cibo ai singoli individui non può essere assicurata secondo i dettami del capitalismo e della privata appropriazione dei frutti del lavoro. Un soldato, anche se mercenario, deve ricevere dal proprio datore di lavoro, quantomeno, armi, vettovaglie e indumenti adeguati. In caso contrario, l’esercito degenererebbe in una congerie di briganti, incapace di sostenere un conflitto bellico moderno. Per condurre una guerra, infatti, un esercito contemporaneo necessita di armamenti tecnologicamente avanzati, di beni di consumo personali indispensabili per preservare la salute e il benessere dei soldati, e di adeguate protezioni individuali. E anche se tutte queste risorse dovessero provenire dal diretto saccheggio delle popolazioni dei territori occupati, si tratterebbe comunque di una necessità logistica, non di una faccenda personale per ciascun soldato, bensì di una condizione imprescindibile per il corretto svolgimento delle loro mansioni.Questa condizione, tuttavia, si distingue nettamente dalle dinamiche del lavoro salariato, fondate sulla natura privata dell'appropriazione. Nelle classiche relazioni capitalistiche di compravendita della forza lavoro, l'appropriazione dei prodotti del lavoro – in tutti i suoi aspetti essenziali e rilevanti per la riproduzione sociale dell'individuo – è separata dal processo produttivo basato sul lavoro retribuito. Il consumo personale è determinato dalla natura privata dell'appropriazione di tali prodotti, mediata attraverso il denaro. E ciò, a sua volta, assicura che il lavoratore sia costretto a riprendere l'attività lavorativa, al fine di procurarsi i beni necessari alla sua sussistenza. In tempo di guerra, la situazione si configura in modo radicalmente diverso: il consumo, in questo contesto, non può essere lasciato, nemmeno in apparenza, alla discrezionalità dei singoli individui.

In definitiva, la guerra moderna rappresenta, tra le altre cose, una competizione tra le diverse forme di produzione non orientate al profitto, sviluppatesi all’interno del sistema capitalista, e cioè la diretta applicazione della forza lavoro umana. In questo scenario, si possono individuare sia l’esercito, inteso come forza militare, sia l’esercito del lavoro, che si organizza in modi differenti. L’attenzione si sposta: se, nel normale funzionamento del capitalismo, ciò che preme è che l’individuo svolga il proprio lavoro al fine di garantirsi la sussistenza, qui, al contrario, l’individuo deve essere in grado di vivere per poter dedicarsi al lavoro. Lo Stato si trova così obbligato ad assumersi questa responsabilità. Inoltre, i rapporti monetari e commerciali diventano un ostacolo evidente al raggiungimento di tali obiettivi.

La guerra sconvolge le consolidate abitudini e le preesistenti forme di attività umana. In tale contesto, lo Stato assume un ruolo preminente, rafforzandosi nell’esercizio delle funzioni distributive all’interno del processo produttivo, sia in ambito militare che nelle retrovie, e si fa carico di stabilire e preservare i legami economici. Tuttavia, in quanto strumento di oppressione di una classe nei confronti di un’altra, o, più precisamente, di una minoranza dominante sulla maggioranza, lo Stato non può risolvere in modo esaustivo tali problematiche, né tantomeno perseguire obiettivi che non siano funzionali agli interessi della classe dirigente. Ciononostante, è proprio durante i conflitti bellici che emergono numerose questioni di questa natura, le quali divergono dagli interessi della classe dominante. Questi problemi si presentano in modo oggettivo, poiché la loro risoluzione si rivela essenziale per la sopravvivenza stessa degli individui.Spesso, lo Stato abdica alla responsabilità di affrontare tali problematiche, relegandole all’iniziativa di svariate organizzazioni e associazioni non governative. Tale inerzia spinge milioni di persone alla ricerca di nuove forme di aggregazione, capaci di farsi carico delle funzioni tradizionalmente svolte dallo Stato. E questi organismi di autogoverno, assumendo compiti che spettano all’ente pubblico, rappresentano, per certi versi, una declinazione alternativa dello Stato stesso.

La questione più urgente e decisiva, la cui risoluzione esige nuove forme di organizzazione in antitesi allo Stato capitalistico, risiede nella prevenzione e nella cessazione delle guerre, nonché nella limitazione dei danni e delle distruzioni che esse arrecano alla società. L'attuazione di questo imperativo cruciale si scontra, tuttavia, con gli interessi del capitale, che non ha ancora esaurito tutte le opportunità di sfruttamento bellico. È fondamentale sottolineare che, in tale contesto, ciò che può assumere una rilevanza storica non è la guerra in sé, né la violenza in quanto tale, bensì la resistenza alla violenza esercitata dalle forze reazionarie, le quali agiscono in nome di interessi egoistici e avvalendosi del sostegno di un'altra classe sociale. È nella resistenza, e nelle sue diverse forme espressive (benché non si limiti ad esse), e non nella guerra in sé, che si manifesta la significativa tendenza all'auto-organizzazione, preludio all'alba della rivoluzione.In assenza di un’adeguata capacità di autorganizzazione, i metodi di produzione umana, diretti e non mediati da beni o denaro, che si sono evoluti all’interno del sistema capitalista, e i meccanismi economici volti alla loro implementazione, non potrebbero costituire il fondamento per la transizione della società verso nuove modalità di sviluppo.

Queste modalità di auto-organizzazione, originate dalla propensione alla guerra e dalla guerra stessa (intesa come reazione alla guerra), e volte a gestire la società, sono per loro natura effimere e destinate a evolvere in tre direzioni possibili: 1) integrazione all’interno dello Stato; 2) degenerazione e successiva dissoluzione (un’integrazione che, tuttavia, non implica necessariamente la degenerazione); 3) assunzione di tutte le funzioni statali, configurandosi come una nuova forma di potere politico. Di conseguenza, tali strutture organizzative rappresentano oggettivamente anche espressioni della lotta di classe, laddove si trovano costrette a contrapporsi allo Stato, assumendo le funzioni di specifici organi statali e privando, in tal modo, l’apparato di potere politico della classe dominante delle sue prerogative.Se in determinate circostanze ciò dovesse effettivamente costituire un beneficio per tali istituzioni, consentendo loro di adempiere alle proprie funzioni in caso di incapacità, allora un ulteriore sviluppo in questa direzione condurrebbe inevitabilmente a uno scontro con lo Stato, a tentativi di limitare le prerogative e l’influenza delle forme di autogoverno popolare. L’esito e la natura di tale conflitto dipenderanno dalle specifiche dinamiche della contesa e dalle strategie adottate dalle parti coinvolte.

Ripercorrendo la storia del Novecento, tra le forme più efficaci di auto-organizzazione che assunsero le prerogative dello Stato, spiccano i soviet. Questi organismi, composti da operai, soldati e contadini, fungevano simultaneamente da organi rappresentativi e legislativi, oltre che da strutture amministrative e di controllo. Sebbene i soviet abbiano fatto la loro comparsa nel 1905, la loro effettiva metamorfosi in una forza statale, in una forma di dualismo di potere, e, soprattutto, l'affermazione dello slogan «Tutto il potere ai soviet», unitamente alle azioni concrete intraprese per tradurre tale principio in realtà, poterono realizzarsi unicamente a seguito del collasso politico del vecchio Stato, innescato da una guerra estenuante. Dopo la Rivoluzione di febbraio, l'inadeguatezza dello Stato, inteso come sistema preposto all'organizzazione della vita sociale, divenne palese, poiché né il governo zarista, né il successivo Governo provvisorio, dimostrarono di possedere gli strumenti per adempiere a tale compito.Mentre i soviet consolidavano la loro posizione, partecipando attivamente alle dinamiche dell’autogoverno locale e dell’amministrazione statale, si delineava uno scenario di dualismo del potere, presente sia nelle retrovie che all’interno dell’esercito. A tal proposito, risulta fondamentale considerare la possibilità che, in futuro, si ripresenti una situazione analoga a quella che caratterizzò la primavera e l’inizio dell’estate del 1917. In quel periodo, si prospettava una transizione pacifica del potere ai soviet, un passaggio fondamentale verso una rivoluzione sociale, che avrebbe creato le premesse per l’avvento del socialismo. Questa possibilità era resa concreta da due fattori: la notevole influenza e il peso sociale dei soviet stessi, e il diffuso armamento della popolazione, conseguenza diretta del conflitto bellico. Lo slogan “Tutto il potere ai soviet” rifletteva, in quel contesto, la consapevolezza dell’esistenza di un simile percorso pacifico verso la rivoluzione. Fu solo a seguito degli eventi di luglio del 1917 che tale via pacifica divenne impraticabile. Ma l’esistenza stessa di questa possibilità rappresenta, per noi, un elemento di primaria importanza.In effetti, qualora tali condizioni si presentassero, è fondamentale non perderle di vista, soprattutto alla luce dell'evoluzione degli armamenti moderni, poiché un approccio pacifico potrebbe rivelarsi l'unica strada praticabile per la trasformazione rivoluzionaria.

Nelle circostanze attuali, possono emergere e stanno emergendo nuove forme di azione, non sempre connesse alla guerra, ma sempre animate dalla necessità di affrontare questioni sociali cruciali. L'essenza di queste modalità di transizione è, ovunque, la stessa: la popolazione si assume direttamente le funzioni del potere statale, soppiantandolo, in tal modo, nella forma a cui siamo abituati nel sistema capitalista.

Passiamo ora a considerare l'organizzazione della vita sociale sotto il comunismo. L'amministrazione della vita sociale si concretizza, in generale, nella pianificazione e nella correzione del piano in base alle reali dinamiche, nell'espressione della volontà collettiva attraverso direttive, nella rendicontazione e nel controllo del processo gestionale, e nell'intervento diretto qualora si renda necessario.

In un sistema comunista, la pianificazione si configura come la pianificazione dello sviluppo umano, divenendo l’obiettivo primario della produzione sociale: la pianificazione della creazione di nuove relazioni interpersonali. Tale processo coincide con la previsione, in quanto la previsione più attendibile è quella per la cui realizzazione vengono compiuti sforzi pianificati. La società deve apprendere a svolgere tutte queste funzioni senza ricorrere a un organo superiore di «gestione», che in realtà si rivela un organo di oppressione, ossia lo Stato.

Con l’avvento di una nuova struttura sociale, ogni funzione legata alla produzione di beni deve essere affidata agli strumenti. Nella fase di transizione verso il comunismo, lo Stato non si limita a dissolversi gradualmente, ma agisce attivamente in tal senso, assicurando che gli individui possano progressivamente farne a meno. Ogni residuo di limitazione della libertà, compresi i meccanismi di gestione del tempo derivanti dall’esistenza dello Stato, si pone in contrasto con il tempo libero, concepito come spazio per lo sviluppo umano e per un’attività libera. Là dove una libertà completa è effettivamente possibile, lo Stato diventa superfluo. Ma laddove questa libertà non sia ancora pienamente realizzata, l’esistenza di una determinata forma di Stato non solo testimonia questa realtà, ma costituisce anche un mezzo per superarla.

Si adduce lo stesso ragionamento contro l’abolizione dello Stato come contro la soppressione della divisione del lavoro. L’essenza di tale argomentazione risiede nel fatto che la vita sociale assume forme sempre più complesse e, di conseguenza, per gestirla adeguatamente, si rende necessaria una differenziazione delle funzioni e un’organizzazione specifica. Tuttavia, da questa crescente complessità non discende affatto l’inevitabilità di corpi amministrativi e di forze armate posti al di sopra della società. Se la società non fosse frammentata in classi, un’associazione di individui autogestita e autonoma sarebbe pienamente possibile, capace di affrontare la complessità della vita sociale con una struttura adeguata, come sostenevano Marx ed Engels, e come ribadiva anche Lenin.

Quando ci riferiamo al comunismo, intendiamo una società nella quale la divisione del lavoro sia stata superata, e in cui gli individui, sviluppando uniformemente le proprie capacità fisiche, la sensibilità, l’intelletto e il senso morale, abbiano acquisito una formazione poliedrica, scientifica, teorica e, al contempo, universale e conforme agli standard moderni e necessari – ovvero al livello più elevato – per il progresso della società. In tal modo, essi sarebbero in grado di apprendere rapidamente qualsiasi mestiere specifico e di impiegare le proprie energie dove si renda necessario, e dove le loro inclinazioni li portino. Ad esempio, un individuo potrebbe dedicarsi prevalentemente allo studio degli uccelli per l’intera esistenza, mentre un altro si dedicherebbe alla danza.

È proprio in virtù di questo principio che, nell'ambito della gestione dello sviluppo sociale mediante l'impiego delle tecnologie disponibili, si persegue un modello comunistico fondato sulla cultura del pensiero sviluppata dall'umanità, in cui tutti sono chiamati a partecipare. Questa rappresenta la premessa fondamentale per il progresso universale di ogni membro della società: per l'intera durata della loro esistenza, tutti devono possedere le competenze necessarie e sviluppare l'abitudine di gestire gli affari pubblici, tutti gli ambiti in cui sono direttamente coinvolti. Una società di questo tipo si rivela più concreta rispetto all'utopico mondo di classe e alle presunte «armoniose relazioni» tra le classi, costantemente propagandate dagli ideologi borghesi contemporanei. Al contrario, il marxismo sottolinea con forza come la mera esistenza dello Stato dimostri l'impossibilità di un simile ordine, nelle condizioni in cui persistono le classi sociali.

Affinché lo Stato cessi di esistere, è necessario che si verifichino determinate condizioni, sia sul piano sociale che su quello tecnico. Per quanto riguarda l’ambito tecnico, un’adeguata formazione politecnica, resa obbligatoria per tutti e fondata sulle più avanzate tecnologie, permetterebbe di organizzare in modo relativamente agevole la gestione dei processi sociali in una società priva di proprietà privata. La condizione più importante, in tal senso, risiederebbe nella formalizzazione della proprietà collettiva dei mezzi di produzione. Naturalmente, per raggiungere tale obiettivo, è imprescindibile superare gli ostacoli di natura sociale, un compito arduo che richiede un impegno considerevole da parte dell’intera società e che non può essere conseguito in tempi brevi. Tuttavia, è importante sottolineare che, dal punto di vista strettamente tecnico, tale obiettivo non appare irraggiungibile già in questa fase. Si tratterebbe di disporre, in tempo reale, di tutte le informazioni relative alla produzione e alla distribuzione in ogni ambito della vita sociale, con la possibilità di tracciare le filiere fino alle loro origini, tenendo conto delle risorse disponibili, delle esigenze manifestate, delle capacità tecnologiche, delle potenzialità comunicative e dei meccanismi di voto, e così via.dall'azione più immediata e diretta fino alla partecipazione ai più ampi processi decisionali collettivi, sia a livello locale che su scala più vasta, il tutto in tempo reale, ogni membro dell'associazione potrà accedere a queste opportunità.

Le modalità di voto e le altre forme di democrazia a cui siamo avvezzi si rivelano efficaci unicamente nelle fasi iniziali di un processo. Man mano che si passa da un sistema di gestione incentrato sulle persone a uno focalizzato su oggetti e processi, l'utilità del voto diminuisce progressivamente, lasciando spazio alla necessità oggettiva intrinseca agli oggetti e alle loro interazioni. Lenin affronta questo concetto in "Stato e rivoluzione", analizzando minuziosamente l'eredità di Marx ed Engels: "Engels parla in modo chiaro e inequivocabile dell''estinzione' e, con espressione ancora più incisiva, del 'sonno' in relazione al periodo successivo alla 'conquista dei mezzi di produzione da parte dello Stato in nome dell'intera società', ovvero dopo la rivoluzione socialista. Sappiamo tutti che, in tale fase, la forma politica dello 'Stato' si manifesta nella sua espressione più compiuta: la democrazia."Ma a nessuno di quegli opportunisti, che manipolano senza scrupoli il marxismo, sfugge il fatto che qui, dunque nelle opere di Engels, si parli dell’“obsolescenza” e del “tramonto” della democrazia. A un primo sguardo, questo potrebbe apparire singolare. Ma è realmente “incomprensibile” solo per coloro che non hanno compreso che la democrazia è anch’essa una forma di Stato e, di conseguenza, che anch’essa è destinata a scomparire con l’estinzione dello Stato stesso. Soltanto una rivoluzione può “abbattere” lo Stato borghese. Lo Stato in generale, ovvero la forma più compiuta di democrazia, può soltanto “decadere”

Nell'epoca contemporanea, segnata da un progresso tecnologico inarrestabile, l'idea di una graduale obsolescenza dello Stato non appare più un'impresa di insormontabile complessità tecnica. In un simile contesto, la gestione e il controllo delle attività e dei processi si riducono a una questione di mera competenza, e i risultati ottenuti diventano patrimonio accessibile a chiunque. Di conseguenza, l'esigenza di una struttura gerarchica si affievolisce, svanendo la premessa che giustificava la separazione di tali funzioni e la loro trasformazione in una professione a sé stante. L'introduzione di sistemi automatizzati per la gestione, unita alla garanzia di un accesso universale alle banche dati e alla possibilità di interazione per tutti gli utenti, rende superfluo il ricorso a una sorta di "esercito di funzionari", i quali, peraltro, potrebbero limitarsi a svolgere un ruolo meramente apparente, configurandosi come un organo distinto e separato dal corpo sociale.

In tal modo, si affronterà una delle sfide più rilevanti nella fase di transizione verso il socialismo: l'istituzione di un sistema di contabilità e controllo, la cui importanza è stata più volte rimarcata da Lenin. Egli affermava che «la contabilità e il controllo rappresentano l'elemento cardine per assicurare il corretto funzionamento della prima fase della società comunista». Si tratta di un sistema di contabilità e controllo applicato alla produzione sociale, e la sua stessa organizzazione costituisce l'embrione di nuove forme di collettività. «Se tutti partecipassero attivamente alla gestione dello Stato, il capitalismo non potrebbe più sopravvivere». Parallelamente, l'evoluzione del capitalismo crea le premesse affinché una partecipazione reale e diffusa alla gestione dello Stato diventi effettivamente possibile.Tra i presupposti fondamentali, si annoverano l’alfabetizzazione generale, già compiuta in molte delle nazioni capitalistiche più avanzate, e l’“istruzione e la disciplina” impartite a milioni di lavoratori attraverso il vasto e articolato sistema collettivizzato delle poste, delle ferrovie, delle grandi industrie, del commercio all’ingrosso, delle banche e di altri enti.

È doveroso ribadire che, tra i presupposti fondamentali, figura ora anche lo sviluppo della tecnologia informatica e le potenzialità in essa racchiuse, le quali, in linea di principio, non possono trovare applicazione in un sistema capitalista, dove manca un'economia pubblica unificata e i rapporti tra i produttori si svolgono attraverso il mercato. Viktor Michajlovič Gluškov ha studiato, nelle sue ricerche, queste potenzialità nell'ambito dell'organizzazione e della gestione della vita sociale, sviluppando il Sistema di gestione economica statale per l'URSS. Gluškov concepiva la gestione della produzione sociale nel suo complesso, avvalendosi di sistemi di controllo che garantissero a tutti l'accesso a una visione dettagliata e in tempo reale della produzione, in modo che le decisioni umane divenissero superflue, poiché il sistema avrebbe agito autonomamente, seguendo schemi predefiniti e conformemente agli obiettivi sociali.

Con il progredire delle tecnologie, risulterà sempre più agevole organizzare la gestione delle risorse e assicurare la partecipazione attiva di tutti i membri della società. Dal punto di vista sostanziale, ciò implica un adeguato sviluppo delle capacità individuali. La vera sfida consiste nel fatto che l’evoluzione verso una società comunista deve iniziare non con coloro che saranno i comunisti del futuro, bensì con gli individui che si sono formati all’interno di un sistema caratterizzato dalla divisione del lavoro, superando le proprie limitazioni attraverso l’istruzione e l’impegno in attività collettive.

L’educazione politecnica si sta affermando come un’esigenza imprescindibile per i membri di un’associazione libera. Il concetto di politecnismo è ormai sufficientemente elaborato, pertanto possiamo rimandare il lettore alle pubblicazioni specialistiche. In questo contesto, tuttavia, desideriamo sottolineare che, al fine di gestire efficacemente i processi sociali, i membri della società devono possedere una visione scientifica globale del mondo e una competenza universale nell’applicazione delle proprie conoscenze. L’odierno funzionario, per adempiere alle proprie mansioni, non necessita affatto di una solida preparazione scientifica, neppure a livello generale; di conseguenza, tale conoscenza tende a scomparire naturalmente, come un elemento superfluo, pur avendo in passato rivestito una certa importanza. Il funzionario rappresenta semplicemente una componente parziale e transitoria di un sistema statale ormai obsoleto.«Una società capace di riorganizzare la produzione secondo principi innovativi, fondati sull’associazione libera e paritaria dei produttori, confinerà l’intera struttura statale al ruolo che le è proprio: un reperto storico, relegato in un museo accanto al fuso e alla celata di bronzo», profetizzava Engels.

Certamente, la gestione dei processi sociali all’interno di un’associazione di questo tipo, e l’impegno attivo nella costruzione della vita sociale, esigeranno molto dai suoi membri, richiedendo una solida capacità di pensiero dialettico, una spiccata sensibilità, integrità morale e un adeguato livello di istruzione. Ma non si tratta di obiettivi irraggiungibili, riservati a un’élite, bensì di qualità che, in determinate condizioni, potranno essere raggiunte da tutti. Man mano che gli individui saranno coinvolti in una vita sociale attiva, partecipando in prima persona, emergeranno come figure di spicco nel panorama culturale. È in quest’ottica che interpretiamo la celebre affermazione di Lenin, secondo cui ogni cuoca dovrebbe imparare a governare lo Stato. L’unico limite a questa possibilità risiede nel modo in cui le persone agiscono oggi, trasformando alcuni individui in semplici funzionari statali e altri in meri esecutori di compiti.

La società comunista è un’associazione libera di individui che organizza il proprio sviluppo sulla base di un piano razionale.

Ma in che modo, precisamente, si autogovernerà? Quali saranno gli strumenti e le misure organizzative che adotterà? A queste domande non possiamo fornire una risposta definitiva. Sarà compito di questa stessa associazione sviluppare, nel corso del tempo, le forme organizzative più adatte alle specifiche condizioni storiche in cui si troverà. L’autogoverno, dunque, si evolverà parallelamente allo sviluppo della società comunista.
\end{document}
